\section{The complete Weil form in logarithmic coordinates}\label{sec:weil}
\subsection{Conventions and the classical criterion}
All Hilbert inner products are conjugate-linear in the first argument.
For $L>0$, let $I_L=(-L/2,L/2)$ and let $E_L$ denote extension by zero
from $I_L$ to $\R$. Unless another domain is specified, $f$ belongs to
$C_c^\infty(I_L;\C)$ and $F=E_Lf$. Our Fourier and translation
conventions are
\begin{equation}\label{eq:fourier}
 \widehat F(\tau)=\int_\R e^{-i\tau x}F(x)\dd x,
 \qquad U_dF(x)=F(x-d),\qquad
 \norm F^2=\frac1{2\pi}\int_\R|\widehat F(\tau)|^2\dd\tau.
\end{equation}
Write $\widetilde F(x)=\overline{F(-x)}$ and $g=F*\widetilde F$.
Then $g(0)=\norm F^2$, $g(-r)=\overline{g(r)}$, and
$g(-r)=\ip F{U_rF}$. In particular, the correlations vanish for $|r|\geq L$.

For $r>0$ set
\begin{equation}\label{eq:gamma-kernel}
 \gammak(r)=\frac{e^{-r/2}}{1-e^{-2r}}
 =\sum_{n\geq0}e^{-a_nr},\qquad a_n=2n+\tfrac12.
\end{equation}
One normalization of the additive Weil functional is
\begin{align}
 \mathcal W(g)={}&\int_\R 2\cosh(x/2)g(x)\dd x
 -\sum_{p,\,m\geq1}(\log p)p^{-m/2}
       \{g(m\log p)+g(-m\log p)\}\notag\\
 &-(\log(4\pi)+\gamma_E)g(0)\notag\\
 &-\int_0^\infty
       [g(r)+g(-r)-2e^{-r/2}g(0)]\gammak(r)\dd r.
 \label{eq:weil-functional}
\end{align}
Here $p$ runs over rational primes and $\gamma_E$ is Euler's constant.
The integrand at the origin has a finite limit, and the prime sum is
finite for compact support. This normalization can be found in
Suzuki's account of the Weil form \cite[Section 1.1]{Suzuki}; for its
geometric and trace interpretation see \cite{CCM}.

The Riemann hypothesis (RH) asserts that every nontrivial zero of the
meromorphic continuation of $\zeta(s)=\sum_{n\geq1}n^{-s}$ has real
part $1/2$. We use the following classical theorem as the number-theoretic
input; we do not reprove the explicit formula or the Weil criterion.
\begin{theorem}[Classical Weil positivity criterion]\label{thm:weil}
RH holds if and only if $\mathcal W(F*\widetilde F)\geq0$ for every
complex $F\in C_c^\infty(\R)$.
\end{theorem}
Every compact support is contained in some $I_L$. Consequently it is
sufficient, and necessary, to establish this positivity for every $L>0$
and every test function in $I_L$.

\subsection{Positive gamma energy and the fixed contact}
Define
\begin{align}
 K[F]&=\frac12\int_\R\int_\R
    |F(x)-F(y)|^2\gammak(|x-y|)\dd x\dd y\notag\\
 &=\int_0^\infty[2g(0)-g(r)-g(-r)]\gammak(r)\dd r.
 \label{eq:gamma-energy}
\end{align}
This is nonnegative without any arithmetic hypothesis. Expanding
\eqref{eq:gamma-kernel} and applying Plancherel gives
\begin{align}
 K[F]&=\frac1{2\pi}\int_\R b(\tau^2)|\widehat F(\tau)|^2\dd\tau,
 \label{eq:gamma-fourier}\\
 b(\tau^2)&=\sum_{n\geq0}\frac2{a_n}
                 \frac{\tau^2}{a_n^2+\tau^2}
 =\re\psi(\tfrac14+i\tau/2)-\psi(\tfrac14),
 \label{eq:gamma-multiplier}
\end{align}
where $\psi=\Gamma'/\Gamma$. Indeed,
$2\int_0^\infty(1-\cos(\tau r))e^{-ar}\dd r
=(2/a)\tau^2/(a^2+\tau^2)$; nonnegativity justifies interchange with
the sum. The second identity in \eqref{eq:gamma-multiplier} is the
digamma partial-fraction expansion \cite{DLMF}.

The archimedean part of \eqref{eq:weil-functional} equals
$K[F]+w_0\norm F^2$, where
\begin{align}
 w_0&= -\log(4\pi)-\gamma_E
       -2\int_0^\infty(1-e^{-r/2})\gammak(r)\dd r\notag\\
 &=\psi(\tfrac14)-\log\pi
 =-\gamma_E-\frac\pi2-3\log2-\log\pi<0.
 \label{eq:contact}
\end{align}
To verify the constant, the integral in the first line is
$\sum_{n\geq0}(a_n^{-1}-(2n+1)^{-1})
=\tfrac12[\psi(1/2)-\psi(1/4)]$.
The constant is therefore fixed by the arithmetic normalization.
It is not a parameter available to adjust a positive source.

\subsection{The signed poles and the localized target}
The first term of \eqref{eq:weil-functional} has kernel
$2\cosh((x-y)/2)$. Using the hyperbolic addition formula gives
\begin{equation}\label{eq:poles}
 P_L[f]=2\left|\int_{I_L}f(x)\cosh(x/2)\dd x\right|^2
       -2\left|\int_{I_L}f(x)\sinh(x/2)\dd x\right|^2.
\end{equation}
It is a rank-two Hermitian form, with one positive and one negative
direction for $L>0$. Its pointwise positive kernel does not make it a
positive quadratic form.

The exact target for this paper is
\begin{align}
 Q_L[f]:=\mathcal W(F*\widetilde F)
 ={}&K[F]+w_0\norm f^2+P_L[f]\notag\\
 &-2\sum_{m\log p<L}(\log p)p^{-m/2}
           \re\ip F{U_{m\log p}F}.
 \label{eq:target}
\end{align}
Its polarization is denoted by $Q_L(f,g)$. At equality $m\log p=L$
the compressed correlation is zero. The coefficient is the primitive
length $\log p$, not the repeated length $m\log p$.

On a fixed interval the prime and pole terms define bounded operators.
The gamma multiplier satisfies
\begin{equation}\label{eq:gamma-growth}
 b(\tau^2)=\log|\tau|-\log2-\psi(\tfrac14)+O(|\tau|^{-2})
 \quad (|\tau|\longrightarrow\infty).
\end{equation}
The natural energy domain is
\begin{equation}\label{eq:energy-domain}
 \mathcal D_{\gamma,L}=
 \left\{f\in L^2(I_L):
   \int_\R\log(2+|\tau|)|\widehat{E_Lf}(\tau)|^2\dd\tau<\infty\right\}.
\end{equation}
The positive gamma energy is closed on this domain; it is the restriction
of the closed nonnegative Fourier-multiplier form to the closed subspace
of functions supported in $\overline I_L$. The bounded other terms give
a closed semibounded form on the same domain. All matching statements
below are first required on smooth compactly supported inputs. The Weil
criterion itself needs only that test space.

Away from the diagonal the distribution kernel of $K$ is
$-\gammak(|x-y|)$, which is smooth whenever $x\ne y$. The prime
contribution instead has atoms at $x-y=\pm m\log p$. This distinction
will provide exact tests for candidate preparations.

Three small intervals already distinguish several issues:
\begin{center}
\begin{tabular}{lll}
\toprule
$L$ & Active prime powers & Test \\
\midrule
$1$ & $2$ & First contact, pole and translation identity \\
$5/4$ & $2,3$ & Joint source and mixed returns \\
$3/2$ & $2,3,4$ & First nontrivial repetition of $2$ \\
\bottomrule
\end{tabular}
\end{center}
A calculation on any one interval is preliminary. It does not establish
the all-support statement of Theorem~\ref{thm:weil}.
